% Chương 5: Kết luận và Hướng phát triển

\chapter{Kết luận và Hướng phát triển}

\section{Kết quả đạt được}

Đồ án đã xây dựng thành công một hệ thống phát hiện xu hướng xã hội hoàn chỉnh, có khả năng tự động thu thập, xử lý và phân tích dữ liệu từ nhiều nền tảng mạng xã hội. Hệ thống kết hợp kiến trúc ba tầng (Producer - Consumer - Dashboard) với các công nghệ hiện đại như Apache Kafka, PySpark, HDBSCAN, và Large Language Models để tạo ra một giải pháp tích hợp và tự động hóa cao.

Về mặt kỹ thuật, hệ thống đã triển khai thành công pipeline xử lý tám giai đoạn với khả năng xử lý 2000-3000 nội dung mỗi lần chạy. Tầng thu thập dữ liệu hoạt động ổn định với ba nguồn TikTok, YouTube và News, sử dụng các API và RSS feeds để lấy dữ liệu mới nhất. Apache Kafka đóng vai trò message queue trung gian, đảm bảo tính ổn định và khả năng mở rộng của hệ thống.

Pipeline xử lý sử dụng PySpark đã tích hợp thành công các kỹ thuật học máy và xử lý ngôn ngữ tự nhiên. Giai đoạn embedding chuyển đổi văn bản thành vector 1536 chiều, cho phép so sánh ngữ nghĩa chính xác. Thuật toán HDBSCAN với cấu hình min\_cluster\_size=15 tự động phát hiện 12-13 cụm xu hướng mỗi batch (trên 2000-3000 items) mà không cần chỉ định trước số lượng cụm. Đặc biệt, việc tích hợp Large Language Models (gpt-4.1-mini, Gemini) cho phép phân tích tự động và trích xuất thông tin chuyên sâu cho từng xu hướng, bao gồm chủ đề, tóm tắt, cảm xúc, từ khóa, loại nội dung, mức độ rủi ro và hướng dẫn sử dụng.

Cơ chế Gộp xu hướng đa yếu tố (embedding similarity, keyword overlap, source distribution) hoạt động hiệu quả, đảm bảo database không bị duplicate và các xu hướng được cập nhật liên tục. Timeout mechanism tự động xóa xu hướng cũ hơn 3 ngày giúp duy trì độ tươi mới của dữ liệu. Checkpoint mechanism ở Consumer đảm bảo exactly-once processing, không bị mất dữ liệu khi xảy ra lỗi.

Về giao diện người dùng, hai dashboard Streamlit cung cấp trải nghiệm tốt cho cả người dùng cuối và quản trị viên. User dashboard hiển thị xu hướng với đầy đủ thông tin, chức năng tìm kiếm và lọc linh hoạt, cùng form đề xuất từ khóa. Admin dashboard cho phép quản lý xu hướng, duyệt đề xuất từ khóa, và theo dõi lịch sử thay đổi. Giao diện trực quan, dễ sử dụng, không yêu cầu kiến thức kỹ thuật.

Từ hai batches thực nghiệm:
\begin{itemize}
    \item Tổng 5305 items được xử lý, lọc ra 4679 items chất lượng cao (88.2\%)
    \item Phát hiện 25 xu hướng từ 2 batches
    \item Noise rate 0.0\% sau two-pass clustering
    \item Accuracy của manual review: 93.3\% (14/15 trends đúng)
    \item Sentiment: 76\% positive, 20\% neutral, 4\% negative
    \item Risk assessment: 76\% safe, 20\% cautious, 4\% high-risk
\end{itemize}

\section{Hạn chế và thách thức}

Bất chấp những kết quả đạt được, hệ thống còn một số hạn chế cần khắc phục trong tương lai.

\textbf{Xử lý thời gian thực:} Hạn chế lớn nhất là chưa hỗ trợ xử lý streaming thời gian thực hoàn toàn. Hệ thống hiện hoạt động theo chế độ batch với chu kỳ bốn giờ, dẫn đến độ trễ trong việc phát hiện xu hướng mới nổi. Mặc dù độ trễ này chấp nhận được cho hầu hết use cases, nhưng với các xu hướng viral nhanh, hệ thống có thể bỏ lỡ giai đoạn bùng nổ ban đầu.

\textbf{Phạm vi dữ liệu hạn chế:} Hệ thống chỉ tập trung vào ba nền tảng TikTok, YouTube và News, chưa bao gồm các nền tảng lớn khác như Facebook, Instagram, Twitter/X. Điều này hạn chế khả năng phát hiện xu hướng đa nền tảng toàn diện. Nguyên nhân chính là hạn chế về thời gian và khả năng truy cập API của các nền tảng này (Facebook và Instagram có API rất hạn chế, Twitter/X tính phí cao).

\textbf{Chi phí vận hành:} Hệ thống phụ thuộc vào các API trả phí như OpenAI (embedding và analysis), Apify (TikTok), và YouTube Data API. Với mỗi batch 2000-3000 items, chi phí embedding và LLM analysis khá cao. Chạy mỗi 4 giờ dẫn đến chi phí hàng tháng đáng kể, có thể là rào cản cho việc triển khai rộng rãi hoặc sử dụng cá nhân.

\textbf{Chất lượng lọc nhiễu:} Quality filtering hiện chỉ dựa trên các chỉ số đơn giản như engagement rate và số followers, chưa phát hiện tốt các loại spam phức tạp hoặc bot tinh vi. HDBSCAN đôi khi tạo ra các cụm quá lớn hoặc quá nhỏ, đặc biệt khi dữ liệu có mật độ không đồng đều. Two-pass clustering đã cải thiện phần nào nhưng vẫn không hoàn hảo.

\textbf{Tính nhất quán của LLM:} LLM analysis đôi khi không nhất quán, đặc biệt với các xu hướng mơ hồ hoặc có nhiều khía cạnh. Hệ thống phụ thuộc hoàn toàn vào chất lượng response từ LLM. Ngoài ra, prompt engineering cần điều chỉnh liên tục để cải thiện chất lượng analysis.

\section{Đánh giá chung}

Nhìn chung, đồ án đã đạt được mục tiêu ban đầu: xây dựng một hệ thống tự động phát hiện và phân tích xu hướng xã hội từ nhiều nguồn. Kiến trúc ba tầng với Kafka làm message queue và PySpark làm processing engine tạo ra một nền tảng vững chắc, có khả năng mở rộng. Việc kết hợp HDBSCAN cho clustering và LLM cho analysis là một approach hiệu quả, tận dụng được ưu điểm của cả machine learning truyền thống và mô hình ngôn ngữ lớn hiện đại.

Hệ thống đã chứng minh khả năng hoạt động trong thực tế, có thể phát hiện các xu hướng thực tế từ dữ liệu mạng xã hội Việt Nam. Dashboard cung cấp giá trị thực tiễn cho các nhà tiếp thị và quản lý mạng xã hội. Những hạn chế hiện tại chủ yếu liên quan đến quy mô và chi phí, có thể được khắc phục trong các phiên bản tiếp theo.

Đồ án không chỉ là một sản phẩm phần mềm hoàn chỉnh mà còn là minh chứng cho khả năng tích hợp nhiều công nghệ tiên tiến vào một hệ thống thống nhất. Kiến thức và kinh nghiệm thu được từ đồ án về xử lý dữ liệu lớn, phân cụm mật độ, embedding vectors, và prompt engineering sẽ là nền tảng vững chắc cho các dự án tương lai.

\section{Hướng phát triển}

Dựa trên kết quả và hạn chế hiện tại, đồ án có thể được mở rộng và cải tiến theo các hướng sau:

\textbf{1. Xử lý streaming thời gian thực:} Chuyển từ batch processing sang streaming architecture sử dụng Spark Structured Streaming hoặc Apache Flink. Điều này cho phép phát hiện xu hướng với độ trễ chỉ vài phút thay vì bốn giờ, đặc biệt quan trọng với các xu hướng viral nhanh. Cần thiết kế lại pipeline để xử lý incremental updates và maintain state cho clustering.

\textbf{2. Mở rộng nguồn dữ liệu:} Tích hợp thêm các nền tảng mạng xã hội lớn như Facebook, Instagram, Twitter/X, và các nền tảng Việt Nam như Zalo, Threads. Mỗi nền tảng có đặc thù riêng về API và loại nội dung, yêu cầu adapter pattern linh hoạt. Ngoài ra, có thể mở rộng sang các nguồn dữ liệu khác như podcast, forum, e-commerce reviews.

\textbf{3. Tối ưu chi phí:} Giảm phụ thuộc vào API trả phí bằng cách:
\begin{enumerate}
    \item Self-host embedding models như sentence-transformers hoặc Gemma
    \item Sử dụng mô hình ngôn ngữ lớn mã nguồn mở như Llama 3 hoặc Qwen cho analysis
    \item Implement caching thông minh để tránh tính toán lại
\end{enumerate}
Trade-off là cần đầu tư infrastructure (GPU) nhưng chi phí dài hạn thấp hơn.

\textbf{4. Cải thiện quality filtering:} Áp dụng các kỹ thuật machine learning để phát hiện spam và bot phức tạp hơn. Huấn luyện classifier dựa trên features như writing patterns, posting frequency, account age. Tích hợp sentiment analysis để lọc nội dung tiêu cực hoặc toxic trước khi phân cụm.

\textbf{5. Dự đoán xu hướng:} Xây dựng mô hình time-series forecasting để dự đoán xu hướng sắp bùng nổ dựa trên tốc độ tăng trưởng engagement, pattern lan truyền, và historical data. Sử dụng kỹ thuật như LSTM, Prophet, hoặc Transformer-based models. Điều này cho phép doanh nghiệp chuẩn bị chiến lược proactive thay vì reactive.

\textbf{6. Personalization:} Xây dựng hệ thống recommendation cá nhân hóa, gợi ý xu hướng phù hợp với từng user dựa trên industry, interests, past interactions. Áp dụng collaborative filtering hoặc content-based filtering kết hợp với embedding similarity.

\textbf{7. Hỗ trợ đa ngôn ngữ:} Mở rộng sang các ngôn ngữ khác ngoài tiếng Việt và tiếng Anh. Sử dụng multilingual embeddings như mBERT hoặc XLM-RoBERTa để xử lý nhiều ngôn ngữ trong cùng một không gian vector, cho phép phát hiện xu hướng cross-language.

\textbf{8. Visual content analysis:} Tích hợp xử lý hình ảnh và video để phân tích visual trends. Sử dụng CLIP hoặc các mô hình multimodal để extract embeddings từ thumbnail, video frames, kết hợp với text embeddings để có phân tích toàn diện hơn.

\textbf{9. Monitoring và alerting:} Xây dựng hệ thống monitoring chi tiết với metrics (processing time, API costs, cluster quality), dashboards cho admin, và alerting khi phát hiện anomalies hoặc xu hướng rủi ro cao. Tích hợp với Prometheus, Grafana hoặc các công cụ APM.

\textbf{10. API và integration:} Cung cấp RESTful API hoặc GraphQL endpoint cho phép các ứng dụng bên ngoài tích hợp dữ liệu xu hướng. Xây dựng webhook để notify real-time khi phát hiện xu hướng mới. Điều này mở ra khả năng commercialize hệ thống như một trend detection service.

